% !TeX program = xelatex
\documentclass[a4paper,10pt]{article}

% ============================================================================
% ក្រដាសដាច់ដោយឡែក៖ សេណារីយ៉ូនៃការបង្រៀន (Scenario នៃការបង្រៀន – ៥០ នាទី)
% ============================================================================
\usepackage[a4paper, left=0.95cm, right=0.95cm, top=0.95cm, bottom=1.05cm, footskip=0.50cm]{geometry}
\usepackage{amsmath}
\usepackage{unicode-math}
\usepackage{fontspec}
\usepackage{array, longtable, colortbl}
\usepackage[table]{xcolor}
\usepackage{pifont}
\usepackage{fancyhdr}
\usepackage{tikz}

% LaTeX 2025 + longtable compatibility guard
\makeatletter
\let\@expl@@@mark@update@singlecol@structures@@\relax
\makeatother

% Cross-platform font path resolver (Windows local + Linux/Cloud repo fonts/)
\providecommand{\NGSFindFonts}{%
  \IfFileExists{C:/Windows/Fonts/cambria.ttc}{%
    \def\NGSSysPath{C:/Windows/Fonts/}%
    \def\NGSUserPath{C:/Users/KIT/AppData/Local/Microsoft/Windows/Fonts/}%
  }{%
    \IfFileExists{./fonts/KhmerOSsiemreap.ttf}{\def\NGSSysPath{./fonts/}\def\NGSUserPath{./fonts/}}{%
    \IfFileExists{../fonts/KhmerOSsiemreap.ttf}{\def\NGSSysPath{../fonts/}\def\NGSUserPath{../fonts/}}{%
    \IfFileExists{../../fonts/KhmerOSsiemreap.ttf}{\def\NGSSysPath{../../fonts/}\def\NGSUserPath{../../fonts/}}{%
    \IfFileExists{../../../fonts/KhmerOSsiemreap.ttf}{\def\NGSSysPath{../../../fonts/}\def\NGSUserPath{../../../fonts/}}{%
    \IfFileExists{../../../../fonts/KhmerOSsiemreap.ttf}{\def\NGSSysPath{../../../../fonts/}\def\NGSUserPath{../../../../fonts/}}{%
    \IfFileExists{../../../../../fonts/KhmerOSsiemreap.ttf}{\def\NGSSysPath{../../../../../fonts/}\def\NGSUserPath{../../../../../fonts/}}{%
    \IfFileExists{../../../../../../fonts/KhmerOSsiemreap.ttf}{\def\NGSSysPath{../../../../../../fonts/}\def\NGSUserPath{../../../../../../fonts/}}{%
      \def\NGSSysPath{../../../../../../../fonts/}\def\NGSUserPath{../../../../../../../fonts/}%
    }}}}}}}%
  }%
}
\NGSFindFonts

\setmathfont[
  Path={\NGSSysPath},
  FontIndex=1,
  Scale=1.02
]{cambria.ttc}

\defaultfontfeatures{Ligatures=TeX}

\setmainfont[
  Path={\NGSUserPath},
  Extension=.ttf,
  Script=Khmer,
  BoldFont={KhmerOSsiemreap},
  BoldFeatures={FakeBold=1.8},
  ItalicFont={KhmerOSsiemreap},
  ItalicFeatures={FakeSlant=0.25},
  BoldItalicFont={KhmerOSsiemreap},
  BoldItalicFeatures={FakeBold=1.8, FakeSlant=0.25}
]{KhmerOSsiemreap}

\newfontfamily\muol[
  Path={\NGSUserPath},
  Extension=.ttf,
  Script=Khmer
]{KhmerOSmuollight}

\newfontfamily\NagaBold[
  Path={\NGSUserPath},
  Extension=.ttf
]{AKbalthom Naga Bold}

\newfontfamily\CalibriFont[
  Path={\NGSSysPath},
  Extension=.ttf
]{calibri}

\XeTeXlinebreaklocale "km"
\XeTeXlinebreakskip = 0pt plus 1pt minus 0.1pt

\definecolor{stepgray}{HTML}{D9D9D9}
\definecolor{headgray}{HTML}{F2F2F2}
\definecolor{darkred}{HTML}{C00000}
\definecolor{red5e}{HTML}{FF0000}

\pagestyle{fancy}
\fancyhf{}
\renewcommand{\headrulewidth}{0pt}
\fancyfoot[R]{\CalibriFont\fontsize{11pt}{13pt}\selectfont\thepage}

\setlength{\parindent}{0pt}
\setlength{\parskip}{0pt}
\setlength{\tabcolsep}{4pt}

\newcolumntype{L}[1]{>{\raggedright\arraybackslash}p{#1}}
\newcommand{\fiveE}[1]{{\NagaBold\fontsize{9pt}{12.5pt}\selectfont\textcolor{red5e}{#1}}}
\newcommand{\ditem}{\tikz[baseline=-0.5ex]\fill (0,0.40ex) -- (0.40ex,0) -- (0,-0.40ex) -- (-0.40ex,0) -- cycle;\hspace{0.25em}}
\newcommand{\chkbox}{\tikz[baseline=-0.1ex]\draw[black, line width=0.75pt](0,0) rectangle (0.23,0.23);}
\newcommand{\cnum}[1]{{\fontsize{9.5pt}{11pt}\selectfont\ding{\numexpr181+#1\relax}}\hspace{0.15em}}
\newcommand{\tri}[1]{\ensuremath{\triangle #1}}
\newcommand{\ang}[1]{\ensuremath{\angle #1}}
\newcommand{\cg}{\ensuremath{\cong}}
\newcommand{\ncg}{\ensuremath{\not\cong}}

\begin{document}
\fontsize{9.4pt}{14.8pt}\selectfont

% -------------------- ក្បាលឯកសារ Scenario --------------------
\begin{center}
  {\muol\fontsize{11pt}{16pt}\selectfont ព្រះរាជាណាចក្រកម្ពុជា}\\[1pt]
  {\muol\fontsize{10pt}{15pt}\selectfont ជាតិ សាសនា ព្រះមហាក្សត្រ}\\[4pt]
  {\muol\fontsize{10.5pt}{15pt}\selectfont សេណារីយ៉ូនៃការបង្រៀន (Scenario នៃការបង្រៀន)}
\end{center}

\vspace{2pt}
{\renewcommand{\arraystretch}{1.06}
\begin{tabular}{@{}l l l p{0.6cm} l l l@{}}
  - អនុវិទ្យាល័យ & : & ព្រះក្រសាំង (កំពង់ស្វាយ-កំពង់ធំ) & & - មុខវិជ្ជា & : & គណិតវិទ្យា ថ្នាក់ទី៨ \\
  - មេរៀនទី១២ & : & ១.២.៣ ជ្រុងបីប៉ុនរៀងគ្នា (ជ.ជ.ជ) & & - រយៈពេលសរុប & : & \textcolor{darkred}{\textbf{១ ម៉ោងសិក្សា (៥០ នាទី)}} \\
  - គ្រូបង្រៀន & : & ហុក ប៊ុនថែន & & - កាលបរិច្ឆេទ & : & ថ្ងៃសុក្រ ទី២៧ ខែវិច្ឆិកា ឆ្នាំ២០២៦ \\
\end{tabular}
}

\vspace{3pt}

% -------------------- I. តារាងបែងចែកពេលវេលា និងសេណារីយ៉ូលម្អិត (៥០ នាទី) --------------------
\textbf{I.} {\muol តារាងបែងចែកពេលវេលា និងសេណារីយ៉ូប្រតិបត្តិការក្នុងថ្នាក់ (៥០ នាទី)}

\vspace{-2pt}
\renewcommand{\arraystretch}{1.12}
\begin{longtable}{|c|L{2.75cm}|c|L{7.65cm}|L{4.45cm}|}
\hline
\rowcolor{headgray}
\multicolumn{1}{|c|}{\muol នាទី} &
\multicolumn{1}{c|}{\muol ជំហាន / 5Es} &
\multicolumn{1}{c|}{\muol រយៈពេល} &
\multicolumn{1}{c|}{\muol សេណារីយ៉ូប្រតិបត្តិ (ពាក្យពេចន៍គ្រូ \& សម្ភារៈ)} &
\multicolumn{1}{c|}{\muol ការគ្រប់គ្រង \& វិធីបម្រុង} \\
\hline
\endfirsthead

\hline
\rowcolor{headgray}
\multicolumn{1}{|c|}{\muol នាទី} &
\multicolumn{1}{c|}{\muol ជំហាន / 5Es} &
\multicolumn{1}{c|}{\muol រយៈពេល} &
\multicolumn{1}{c|}{\muol សេណារីយ៉ូប្រតិបត្តិ (ពាក្យពេចន៍គ្រូ \& សម្ភារៈ)} &
\multicolumn{1}{c|}{\muol ការគ្រប់គ្រង \& វិធីបម្រុង} \\
\hline
\endhead

\hline
\endfoot

\hline
\endlastfoot

00--02 & \textbf{ជំហានទី១}\newline រដ្ឋបាលថ្នាក់ & ២ នាទី &
\ditem ហៅអវត្តមានរហ័ស និងពិនិត្យអនាម័យថ្នាក់រៀន ។\newline
\ditem ចាត់សិស្សជាក្រុម ៣ នាក់ («ក្រុមជាងសំណង់ដំបូល») ដាក់កញ្ចប់កាត់ ១:១ (ទំព័រ ១) ទុកមុខតុ។ &
ចាត់ប្រធានក្រុមស្រាប់ពីមុន ហាមបើកកញ្ចប់កាត់មុនបញ្ជា ។ \\ \hline

02--05 & \textbf{ជំហានទី២ (ក)}\newline កំណែកិច្ចការផ្ទះ & ៣ នាទី &
\ditem \textbf{ល្បែងចាប់កំហុស}៖ សួរបញ្ឆោត «តើអាចយក $AB=AC$ មកដាក់ក្នុង $\tri{BDE}, \tri{CDE}$ បានឬទេ?»\newline
\ditem ផ្ទៀងចម្លើយ SAS ទាំង ក ($DB=DC, \ang{ADB}=\ang{ADC}$) និង ខ ($\tri{BDE}\cg\tri{CDE} \Rightarrow EB=EC$)។ &
កុំហៅសរសេរឡើងវិញលើក្ដារខៀន ផ្ដោតតែការបញ្ជូន $DB=DC$ ពី ក ទៅ ខ ។ \\ \hline

05--08 & \textbf{ជំហានទី២ (ខ)}\newline រំលឹកមេរៀនទាក់ទង & ៣ នាទី &
\ditem \textbf{ភាគ ១ (១.៥ ន.)}៖ គូសធ្នូកាំ $5\text{cm}$ ផ្ចិត $A$ បញ្ជាក់គ្រប់ចំណុចលើធ្នូមានចម្ងាយ $5\text{cm}$ ពី $A$ ។\newline
\ditem \textbf{ភាគ ២ (១.៥ ន.)}៖ បង្ហាញ $F\text{-}J\text{-}G\text{-}K$ រត់ត្រង់ជួរ $FJ=GK \Rightarrow FG = FJ+JG = GK+JG = JK$ ។ &
បើសិស្សភ្លេចបូក $JG$ ឱ្យម្រាមដៃអូសពី $F\to G$ និង $J\to K$ ភ្លាម ។ \\ \hline

08--11 & \fiveE{E1 Engage}\newline អាថ៌កំបាំងដំបូល & ៣ នាទី &
\ditem លើករឿងពិត៖ «ជាង ២ នាក់នៅកំពង់ស្វាយកាត់ដំបូល $6\text{m}, 5\text{m}, 3\text{m}$ ដូចគ្នាដោយ\textbf{គ្មានរ៉ាព័រទ័រ} ហេតុអ្វីរាង/មុំថេរ ខណៈទ្វារ ៤ របារទ្រេតដួល?»\newline
\ditem ស្រង់មតិទស្សន៍ទាយ រួចសរសេរចំណងជើង ១.២.៣ ។ &
កុំទាន់លាតត្រដាងចម្លើយ ទុកឱ្យសិស្សចង់សាកល្បងនៅ E2 ។ \\ \hline

11--22 & \fiveE{E2 Explore}\newline ពិសោធន៍ ១ \& ២ & ១១ នាទី &
\ditem \textbf{សកម្មភាព ១៖ ស៊ុម ៤ vs ៣ របារ \& ប្រសព្វធ្នូ (៧ ន.)}៖\newline
+ ២.៥ន. ចាក់ម្ជុល ៤ របារ ($6,5,4,3\text{cm}$) តេស្តការទ្រេតប្រែរាង\newline
+ ២.៥ន. តភ្ជាប់ $AB=6, AC=5, BC=3\text{cm}$ + ធ្នូ $r_1=5, r_2=3\text{cm}$ ប្រសព្វ $C$ តែ \textbf{១ គត់}\newline
+ ២ន. ត្រួតបង្វិល-ត្រឡប់ក្រដាសកាត់លើក្រុមជិតខាង ១០០\% ។\newline
\ditem \textbf{សកម្មភាព ២៖ POE អន្ទាក់ ម.ម.ម (៤ ន.)}៖ ៣០វិ. ទស្សន៍ទាយ (P) $\rightarrow$ ២.៥ន. ត្រួត $\tri{PQR}$ ($4\text{cm}$) លើ $\tri{P'Q'R'}$ ($6\text{cm}$) មុំ $60^\circ$ ច្រាន ម.ម.ម (O–E) + ១ន. តេស្ត $2, 2.5, 6\text{cm}$ ។ &
បោះពុម្ពសម្ភារៈទំព័រ ១ តាមខ្នាត \textbf{Actual Size 100\%} ។ គ្រូត្រួតការប្រើម្ជុល និងដែកឈានប្រសព្វ $C$ ។ \\ \hline

22--28 & \fiveE{E3 Explain}\newline ទ្រឹស្ដីបទ \& កំពូល & ៦ នាទី &
\ditem បង្គុំទ្រឹស្ដីបទ ជ.ជ.ជ លើរូបបង្វិល-ឆ្លុះ $\approx 135^\circ$ + \textbf{គោលការណ៍កំពូលឈមជ្រុងស្មើ} ។\newline
\ditem ដឹកនាំរកកំហុសបោះពុម្ពទំព័រ ១៥២ (\textbf{$AB=B'C' \rightarrow BC=B'C'$}) បិទផ្ទាំង \textbf{ធនាគារហេតុផល \cnum{1}--\cnum{5}} (ទំព័រ ៤) និងញែក $=$ vs $\cg$ ។ &
សង្កត់ធ្ងន់៖ ដាក់រូបបង្វិល-ឆ្លុះ កុំគូសតែទម្រង់ត្រង់ស្ដង់ដារ ។ \\ \hline

28--38 & \fiveE{E4 Elaborate}\newline ស្ពាន $JG$ (\textbf{[tracing]}) & ១០ នាទី &
\ditem ចែកកាត \textbf{«ស្ពានតភ្ជាប់អង្កត់ $JG$ និងគំនូស $|, \parallel, |||$»} (ទំព័រ ២, ទំព័រគោល ១៥២)៖\newline
+ ៣ន. ប្រើម្រាមដៃ \textbf{[tracing]} គូស $|$ ($EF=HK$), $\parallel$ ($FJ=GK$), $|||$ ($EG=HJ$) និងអូសស្ពាន $JG$\newline
+ ៥.៥ន. សរសេរភស្តុតាង \cnum{1}--\cnum{4} ($\tri{EFG}\cg\tri{HKJ}$) + វិបាក \cnum{5} ($\ang{E}=\ang{H}$)\newline
+ ១.៥ន. \textbf{EIA ពិនិត្យក្រោយ (30s)} រំកិល $\tri{HKJ}$ រក្សា $FJ=GK$ ញែកប្រែ ($JG,FG,JK$) vs ថេរ ($FG=JK$) ។ &
ដាក់លេខ \cnum{2} នៅជួរ $FG=JK$ កុំឱ្យសិស្សលោតរំលង $JG$ ។ គូស $|||$ លើផ្នែកខាងលើ $EG, HJ$ កុំឱ្យប៉ះចំណុចប្រសព្វ ។ \\ \hline

38--48 & \fiveE{E5 Evaluate}\newline អធិការកិច្ច \& កាត A--D & ១០ នាទី &
\ditem \textbf{ក្រុមអធិការកិច្ចប្លង់ ៣ តួនាទី (៧ ន.)} (ទំព័រ ៣, ប្រតិបត្តិទំ.១៥៣)៖\newline
+ \textbf{អ្នកទី១ (Pause-Ask)}៖ ពន្យល់មុនគូស $|, \parallel, |||$ (\textbf{ហាមគូស $RT=SV$})\newline
+ \textbf{អ្នកទី២}៖ អូសស្ពានឈរ $ST$ ស្រាយ $RT=SV$ (\cnum{2}) និង \cnum{1}--\cnum{4}\newline
+ \textbf{អ្នកទី៣ (Pause-Review)}៖ ផ្ទៀង $R\leftrightarrow V, Z\leftrightarrow W, T\leftrightarrow S$ ទាញ \cnum{5} \& \textbf{បិទ E1} ។\newline
\ditem \textbf{កាតត្រួតពិនិត្យ A--D (៣ ន. / 60--90s)} (ទំព័រ ៥)៖ (A) \ding{51} ជ.ជ.ជ $135^\circ$, (B) \ding{55} ម.ម.ម $60^\circ$, (C) \ding{55} $FJ=GK$ ខ្វះស្ពាន, (D) \ding{55} លំដាប់កំពូលខុស + ម្រាម ១--៣ ។ &
បើ E2 ចំណាយលើស ១ នាទី កាត់កាត A--D មក ៩០ វិនាទីរហ័ស ដើម្បីបញ្ចប់ទាន់ ៥០ នាទី ។ \\ \hline

48--50 & \textbf{ជំហានទី៥}\newline កិច្ចការផ្ទះ \& បណ្ដាំ & ២ នាទី &
\ditem ចែកប័ណ្ណកិច្ចការផ្ទះ (ទំព័រ ៦)៖ លំហាត់ទី៤ (ទំ.១៦៤ ខ្លែង $ABCD$) \& ទី៩ (ទំ.១៦៥ ប្រយ័ត្ន $\tri{ABC}\cg\tri{DCB}$ កុំច្រឡំ $BDC$) ។\newline
\ditem ផ្ដាំផ្ញើសុវត្ថិភាពចរាចរណ៍ និងជួយការងារឪពុកម្ដាយ ។ &
បញ្ជូនរូបលំហាត់ចូល Telegram / Classroom ជាស្រេច ។ \\ \hline

\rowcolor{stepgray}
\multicolumn{2}{|r|}{\textbf{សរុបរយៈពេលបង្រៀនពេញ}} & \textbf{៥០ នាទី} &
\multicolumn{2}{l|}{ត្រូវគ្នា ១០០\% នឹង \texttt{LS\_Plan\_SSS.tex} និង \texttt{Printable\_Materials\_SSS.tex} (ទំព័រ ១--៦)} \\ \hline
\end{longtable}

\vspace{2pt}

% -------------------- II. បញ្ជីត្រៀមសម្ភារៈ & III. ការគ្រប់គ្រងកំហុស --------------------
\begin{minipage}[t]{0.46\textwidth}
  \textbf{II.} {\muol បញ្ជីត្រៀមសម្ភារៈមុនចូលថ្នាក់ (ស-ខ្មៅ)}\\[3pt]
  \begin{tabular}{|c|L{7.2cm}|}
    \hline
    \rowcolor{headgray}
    \textbf{\ding{51}} & \multicolumn{1}{c|}{\textbf{សម្ភារៈត្រូវរៀបចំទុកជាមុន}} \\ \hline
    \chkbox & សន្លឹកកាត់ ១:១ ទំព័រ ១ (\textbf{Actual Size 100\%})៖ របារ $6,5,4,3\text{cm}$, របារ $2,2.5,6\text{cm}$ \& ត្រីកោណ AAA $4,6\text{cm}$ \\ \hline
    \chkbox & ម្ជុលខ្ទាស់/ប៊ូតុងចោះរន្ធ ដែកឈាន និងបន្ទាត់ក្រិតសម្រាប់សង់ធ្នូប្រសព្វ $C$ \\ \hline
    \chkbox & កាតសកម្មភាព E4 (ទំព័រ ២), សន្លឹកអធិការកិច្ច ៣ តួនាទី E5 (ទំព័រ ៣) \& ផ្ទាំងធនាគារ \cnum{1}--\cnum{5} (ទំព័រ ៤) \\ \hline
    \chkbox & កាត Exit \textbf{A--D} (ទំព័រ ៥) \& ប័ណ្ណកិច្ចការផ្ទះទី៤, ទី៩ (ទំព័រ ៦) \\ \hline
  \end{tabular}
\end{minipage}%
\hfill
\begin{minipage}[t]{0.51\textwidth}
  \textbf{III.} {\muol កំហុសភ័ន្តច្រឡំ (តាមការស្រាវជ្រាវ) \& ដំណោះស្រាយ}\\[3pt]
  \begin{tabular}{|L{8.8cm}|}
    \hline
    \rowcolor{headgray}
    \multicolumn{1}{|c|}{\textbf{ឧបសគ្គសិស្ស (Empirical Misconceptions) \& វិធានការទប់ស្កាត់}} \\ \hline
    \textbf{១. អន្ទាក់ ម.ម.ម (AAA) \& រូបបង្វិល-ឆ្លុះ (ធ្លាក់ដល់ ២១.៥\%)៖} ត្រួតផ្ទាល់ $\tri{PQR}$ ($4\text{cm}$) លើ $\tri{P'Q'R'}$ ($6\text{cm}$) និងប្រើ \textit{គោលការណ៍កំពូលឈមជ្រុងស្មើ} លើរូបឆ្លុះ $135^\circ$ ។ \\ \hline
    \textbf{២. រូបត្រួតគ្នា (ធ្លាក់ពី ៨៤\%$\to$៤៩\%) \& រំលងស្ពាន $JG/ST$៖} បង្ខំ \textbf{[tracing]} អូស $JG/ST$ និង \textbf{Pause-Ask} ហាមគូស $RT=SV$ មុនសរសេរ \cnum{2} ។ \\ \hline
    \textbf{៣. ភ័ន្តច្រឡំធ្នូដែកឈាន \& $=$ vs $\cg$ (៦៤\%)៖} បញ្ជាក់ធ្នូកាំ $5,3\text{cm}$ ជាចម្ងាយថេរ និងប្រើផ្ទាំង \cnum{1}--\cnum{5} ញែក $=$ (ប្រវែង/មុំ) ពី $\cg$ (ត្រីកោណ) ។ \\ \hline
    \textbf{៤. កំហុសសៀវភៅទំ.១៥២ \& លំហាត់ទី៩៖} កែ $AB=B'C'\to BC=B'C'$ ក្នុង E3 និងផ្ទៀង $\tri{ABC}\cg\tri{DCB}$ ជំនួស $BDC$ ។ \\ \hline
  \end{tabular}
\end{minipage}

\end{document}
